\section{Examples of cover assignment schemes} \label{sec:examples}

We now give example strategies to define cover assignment schemes, beginning with the one that corresponds to the standard Mapper construction defined in Section~\ref{sec:background}.


\subsection{Standard cover assignment scheme}
\label{sec:stdMapper}


%It is possible to link our construction to the classical Mapper graph construction, i.e. the Mapper graph built to approximate the Reeb graph of a continuous function.\\
%\noindent As before, $\Xs$ is a point cloud of size $n$ in  a general metric space$(\Metrsp,d)$. As explained before, a deterministic clustering algorithm has been fixed for defining the function $\MapComp$. 
Let $f\colon\Xs_n \rightarrow \R$ be a filter function  and let $(\interv_j)_{1\leq j\leq r}$ be a finite cover of the image $f(\Xs_n)$ of $f$.
%Let $r'=\sum_{j=1}^r k_j$, where $k_j$ is the number of clusters detected by the fixed clustering algorithm $\Clus$ on $f^{-1}(I_j)$. 
The standard Mapper graph is then defined as $\MapComp(e^\ast)$, where for every $1\leq i \leq n$ and 
%$1\leq j'=(j,k) \leq r'$ (with $1\leq j\leq r$ and $1\leq k\leq k_j$):
$1\leq j \leq r$:
$$e^\ast_{i,j}=1 \text{ if } f(x_i)\in \interv_j.$$ %\text{ and }x_i\in \mathcal C_{j,k}.$$
The cover assignment scheme $A^\ast$, in this case, is such that every entry $A^\ast_{i,j}$ follows a Dirac distribution on $1$ if $f (x_i)\in \interv_j$,
%and $x_i\in \mathcal C_{j,k}$, 
and a Dirac distribution on  $0$ otherwise. In other words, the parameters of the Bernoulli distributions satisfy $p_{i,j} (\Xs_n) = p_{i,j} (x_i) = 1 $ if  $f (x_i)\in \interv_j$,
%and $x_i\in \mathcal C_{j,k}$, 
and $0$ otherwise, that is
\[
    \mathbb{P}(A^\ast=e | \Xs_n)= 
\begin{cases}
    1& \text{ if } e=e^\ast ,\\
    0              & \text{ otherwise.}
\end{cases}
\]
In this degenerated situation, the random variables $A^\ast_{i,j}$ are all independent conditionally to $\Xs_n$, and $A^\ast_{i,j}$ conditionally to $\Xs_n$ is equal to $A^\ast_{i,j}$ conditionally to $x_i$. 



%Note also that the different coordinates of $A^\ast$ conditionally follow Dirac distributions in $\{0,1\}$, but we assume that they can be seen as an extension of the Bernoulli distribution when the success probability is either 1 or 0.\\
% The conditional risk of the cover assignment scheme $A^\ast$ is : 
% $$\mathbb{E}_{\mathbb{X}_n,\theta}(\mathcal{L}(A^\ast,\theta))=\mathcal{L}(e^\ast,\theta).$$

\begin{remark}
    An alternative and relevant approach for the standard Mapper graph is to define the intervals $I_j$ via the quantiles of the distribution of $f(\Xs_n)$.  
In this case, the random variables $A^\ast_{i,j}$ do not only depend on $x_i$, but also on the whole point cloud $\Xs_n$.
\end{remark}


\subsection{Smooth relaxation of the standard cover assignment scheme}
\label{sec:Smooth}

Given some $\delta>0$, we can now define a cover assignment scheme $A_\delta$ that approximates the cover assignment scheme $A^\ast$ arising from the standard Mapper graph,
but that also enjoys useful smoothness properties in the optimization setting that we will consider in the next section. %, when $\delta$ is close to zero.\\
Specifically, using the same notations as before, and denoting each element of the cover with $I_j=[a_j,b_j]$, consider, for each $j\in\{1,...,r\}$, the function $q_j\colon X  \longrightarrow [0,1]$ defined with:
\begin{align*}
&x \mapsto 
\begin{cases}
    1,& \text{if } f(x)\in [a_j,b_j]\\
    \exp(1-1/(1-(\frac{a_i-f(x)}{\delta})^2)),& \text{if } f(x)\in [a_j-\delta,a_j]\\
    \exp(1-1/(1-(\frac{f(x)-b_i}{\delta})^2)),& \text{if } f(x)\in [b_j,b_j+\delta]\\
    0,              & \text{otherwise}
\end{cases}
\end{align*}
\noindent
Now, define $A_\delta=(A_{\delta,i,j})_{\substack{1\leq i\leq n \\ 1\leq j\leq r}}$ to be the random variable in $\{0,1\}^{n\times r}$ such that for every $(i,j)\in\{1,...,n\}\times\{1,...,r\}$:
$$ A_{\delta,i,j}\mid \Xs_n \sim \mathcal{B}(q_j(x_i)),$$
%and we take 
with the $ A_{\delta,i,j}$'s being jointly independent conditionally to $\Xs_n$.
As for the standard cover, the Bernoulli parameter $p_{i,j} = q_j(x_i)$ depends on its associated point $x_i$ and also on the chosen filter $f$. 
%This remark will be important in Section~\ref{sec:optim}.


Moreover, notice that for every $x_i\in\Xs_n$ and $j\in\{1,...,r\}$:

$$
    q_j(x_i)\xrightarrow[\delta \to 0]{}\begin{cases}
    1,& \text{if } f(x_i)\in I_j\\
    0,              & \text{otherwise,}
\end{cases}
$$
and this shows that $A_\delta\xrightarrow[\delta \to 0]{\mathcal{L}}A^\ast.$ Note that even though we can approximate the standard Mapper graph in this way, we do not always want to do so. For example, there could be cases where $\delta$ needs to be large enough so as to account for some uncertainty on the bounds of the cover $(\interv_j)_{1\leq j\leq r}$. 


%The reason behind considering such an approximation of $A^\ast$ is because it allows for useful smoothness properties in the optimization setting that we will consider in the following section.  
\begin{remark}
Note that the same relaxed construction can be made for a multi-dimensional Mapper, i.e., for filter functions taking values in $\R^d$ \cite{carriere2022statistical}, by making slight adjustments to the definition of $q_j$ %, for every $j\in\{1,...,r\}$, 
using the Euclidean norm.
\end{remark}
 


An additional example of a possible cover assignment scheme, which does not imply the existence of a filter function, is given in Appendix \ref{apdx:gaussian}.

